% ==============================================================================
% SECCIÓN 7: GESTIÓN ÁGIL CON TRELLO (HERRAMIENTA SELECCIONADA)
% ==============================================================================

\section{Gestión Ágil del Proyecto con Trello}

En cumplimiento de las pautas del laboratorio y adoptando un marco de trabajo visual ágil (Kanban), se ha diseñado la estructura operativa del proyecto en la plataforma \textbf{Trello}. Esta configuración permite una trazabilidad completa de los entregables de la fase de \textit{Inception} y de las futuras historias de usuario del sistema.

\subsection{Estructura del Tablero Kanban (SIGC\_V0.1)}

El tablero se organiza en cinco columnas principales con políticas de flujo explícitas:

\begin{enumerate}
    \item \textbf{01. Product Backlog (Épicas del Sistema):}
    Contiene los requerimientos macro del sistema clasificados por módulos:
    \begin{itemize}
        \item \textit{Épica 1:} Gestión Multientidad (Municipalidades y Empresas Turísticas).
        \item \textit{Épica 2:} Catálogo de Capacitaciones y Motor de Inscripción Pública.
        \item \textit{Épica 3:} Control de Asistencia Digital en Tiempo Real (Marcado QR).
        \item \textit{Épica 4:} Registro Académico y Actas de Evaluación.
        \item \textit{Épica 5:} Emisión Automatizada y Validación Pública de Certificados QR.
        \item \textit{Épica 6:} Dashboard Gerencial y Reportes de Impacto.
    \end{itemize}

    \item \textbf{02. Sprint Inception (To Do):}
    Actividades inmediatas correspondientes al Laboratorio 02:
    \begin{itemize}
        \item Redactar Documento de Visión bajo la plantilla de Geoffrey Moore.
        \item Diseñar y validar el Lean Canvas completo en formato horizontal.
        \item Elaborar la Matriz de Stakeholders con necesidades y dolores de Cusco.
        \item Identificar riesgos tempranos y formular planes de contingencia técnica.
        \item Configurar repositorio colaborativo en GitHub con políticas de ramas.
    \end{itemize}

    \item \textbf{03. In Progress (En Proceso):}
    Tareas asignadas a los integrantes del equipo (Choquenaira, Yaranga, Ccama) con un límite estricto de trabajo en curso (\textit{WIP Limit} = 3 tarjetas activas simultáneamente para evitar cuellos de botella).

    \item \textbf{04. Review / QA (Revisión de Calidad):}
    Revisión de la documentación y artefactos técnicos: verificación de formato APA 7ma edición, compilación limpia en \LaTeX{} y comprobación de criterios de aceptación.

    \item \textbf{05. Done (Completado y Validado):}
    Entregables listos para su presentación académica y posterior paso a la fase de diseño arquitectónico.
\end{enumerate}

\subsection{Esquema de Etiquetas y Asignaciones}
Para una identificación visual inmediata en Trello, se establecen etiquetas cromáticas estándar:
\begin{itemize}
    \item \textbf{Verde:} Módulo de Certificación y Validación QR.
    \item \textbf{Azul:} Módulo de Capacitaciones e Inscripción.
    \item \textbf{Naranja:} Control de Asistencia y Evaluaciones.
    \item \textbf{Rojo:} Gestión de Riesgos y Seguridad Criptográfica.
    \item \textbf{Púrpura:} Documentación y Metodología Agile Inception.
\end{itemize}

\subsection{Definición de Terminado (\textit{Definition of Done - DoD})}
Una tarjeta en Trello solo puede promoverse a la columna \textbf{Done} si cumple concurrentemente con:
\begin{enumerate}
    \item Redacción técnica rigurosa y sin ambigüedades.
    \item Integración verificada en el documento formal \LaTeX{} sin advertencias críticas de compilación.
    \item Revisión y aprobación por al menos dos miembros del equipo mediante un \textit{Pull Request} en GitHub.
\end{enumerate}
